\section{RL as probabilistic inference: una interpretación útil, no una demostración universal}

A continuación, el ponente sitúa Decision Transformer dentro de la perspectiva de «control as inference»: interpreta una reward alta como una mayor optimality probability de la trayectoria y luego selecciona las acciones mediante inferencia condicional. Esta perspectiva ayuda a entender el return conditioning, pero no sustituye la cobertura de los datos fuera de línea, la identificación causal ni la evaluación en el entorno real.

\subsection{La intuición de la optimality variable}

Se introduce una optimality variable binaria $\mathcal{O}_t$ y se define a partir de la reward
\[
p(\mathcal{O}_t=1\mid s_t,a_t)\propto \exp\left(\frac{r(s_t,a_t)}{\alpha}\right),
\]
donde $\alpha>0$ es la temperature. Los pares state-action con reward alta tienen más probabilidad de considerarse optimal. La distribución posterior de la trayectoria completa condicionada a ``optimal'' puede escribirse como la combinación de la prior dynamics y la reward weighting; el policy learning puede entenderse como una aproximación de esa distribución posterior.

\lecturefigure{slide-13-rl-as-probabilistic-inference.jpg}{RL as probabilistic inference: las optimality variables convierten la reward en un problema de inferencia condicional.}{Video oficial de Stanford Online, 30:00; la figura de la derecha cita a Levine 2018.}

\begin{knowledgebox}{Lectura de la figura: la conexión entre Decision Transformer y la inference view}
A la izquierda sigue estando la token sequence $(R,s,a)$; a la derecha, el probabilistic graphical model conecta la optimality variable de cada paso con el state/action. La conexión que plantea la clase es la siguiente: el return-to-go puede verse como una condición comprimida sobre la optimality futura, y el Transformer aprende la action distribution dada esa condición. No ejecuta con exactitud la inferencia del modelo gráfico, sino que aproxima con datos supervisados la distribución condicional correspondiente.
\end{knowledgebox}

\subsection{Qué explica esta perspectiva}

Explica por qué «condicionar a un retorno alto» puede inclinar la elección hacia acciones de mayor calidad, y también por qué el behavior prior es importante: si cierta acción nunca aparece en los registros, a la posterior approximation le resulta difícil respaldarla de forma confiable. Además, sugiere la relación entre la temperature, la reward scale y la entropy regularization, lo que ayuda a conectar con el maximum-entropy RL.

\subsection{Qué no demuestra esta perspectiva}

La interpretación probabilística no demuestra por sí sola que el return objetivo sea alcanzable, ni que el modelo condicional extrapole correctamente fuera del dataset support; mucho menos demuestra que el Transformer ya realice planning causal. Lo que el modelo aprende es la estructura de correlaciones en la observational trajectory distribution; la dinámica del entorno y la action consequence todavía tienen que estar reflejadas en los datos.

\begin{warningbox}{Un patrón de trayectorias correlacionado no es una garantía contrafactual}
Si las trayectorias con return alto siempre van acompañadas de cierto rasgo visual sin efecto causal, el modelo puede tomarlo como una señal para actuar. Solo con intervention, aleatoriedad del entorno, datos de distintas policies o una counterfactual evaluation específica es posible distinguir con más certeza la verdadera action consequence de las correlaciones presentes en los registros.
\end{warningbox}

\subsection{Resumen de la sección}

RL-as-inference ofrece una intuición unificada para el return conditioning: las trayectorias con reward alta reciben una optimality weight mayor y la policy aproxima la distribución posterior condicional. Sin embargo, esta interpretación sigue sujeta al behavior prior y al soporte de los datos, y no sustituye la verificación empírica en lazo cerrado.
